\section{Methodology}\label{sec:methodology}

\subsection{Overview}

AutoEvolve starts from $\mathcal{H}_{\min}$ and evolves toward $\mathcal{H}_{\text{auto}}$ during a single continuous episode. The agent begins with only frame observations, button inputs, a generic system prompt, and meta-tools for creating harness components. It receives no pathfinding, no walkthrough, and no external knowledge. Everything about the game must be discovered through play.

Two loops run in interleaved fashion. The inner loop is the standard agent step cycle: observe, reason, act. The outer loop is evolutionary optimization: every $F$ steps (we use $F{=}50$, after a warmup of $W{=}50$ steps), a meta-model analyzes recent trajectory segments and refines all four evolvable components: the system prompt $p$, the sub-agent registry $\mathcal{S}$, the skill library $\mathcal{K}$, and the memory store $\mathcal{M}$. The agent plays without interruption; evolution results are injected into the orchestrator's context on the next step.

The meta-tools let the agent create and manage its own harness during gameplay. It can CRUD sub-agents in a persistent registry (each with a custom prompt, tool whitelist, and termination condition), manage a skill library where entries can be text descriptions or executable Python, maintain a hierarchical knowledge store, launch sub-agents that run autonomously with fresh observations each turn, and analyze any window of past trajectory steps.

\seth{@Joel TODO: GPP meta-tool mapping. Which GPP tools map to which AutoEvolve components? Specific examples: press\_sequence, truth-table notepad, Operation Zombie Phoenix.}

\subsection{Reset-Free Evolutionary Optimization}\label{sec:evolution}

GEPA~\citep{agrawal2025gepa} evolves only $p$ by running complete episodes, scoring them, rewriting the prompt, and resetting. This requires $K$ full rollouts for $K$ evolution steps and cannot adapt to late-game situations.

AutoEvolve operates on trajectory segments instead. The HarnessEvolver executes four independent passes per cycle. \textbf{Prompt evolution}: the meta-model rewrites the strategic guidance given the last $F$ trajectory entries and the fixed system prompt (tool docs, constraints) as read-only context. \textbf{Sub-agent evolution}: given the registry and recent trajectories (including detected tool failures), the meta-model recommends new sub-agents for repeated manual patterns, updates to existing ones, or retirement of ineffective ones. \textbf{Skill evolution}: the meta-model identifies reusable patterns from successful sequences, codifies them as skills, updates effectiveness ratings, and flags broken executable code for repair. \textbf{Memory curation}: a lightweight pass that fills knowledge gaps, updates stale entries, and rebalances importance scores.

\subsection{Sub-Agents and Skills}\label{sec:subagents-skills}

Sub-agents are autonomous modules the orchestrator delegates to. Each registry entry specifies a system prompt (up to 12K characters), a tool whitelist, a handler type (one-step analysis or multi-turn loop), a termination condition, and a step budget. When invoked, a sub-agent loops with fresh screenshots and game state each turn until it signals completion or exhausts its budget. The orchestrator receives a compacted summary. Sub-agents cannot spawn other sub-agents.

Sub-agents emerge through two mechanisms. The orchestrator creates them when it identifies a recurring need (for instance, entering a battle). The evolver creates them when it detects repeated multi-step manual patterns in trajectories. In practice both happen: agents create their first sub-agents within 20 to 30 steps, and the evolver refines them at each cycle.

Skills take two forms. Behavioral skills are text descriptions referenced in reasoning. Executable skills are Python code running in a sandbox with access to \texttt{press\_buttons}, \texttt{get\_game\_state}, and standard libraries (\texttt{numpy}, \texttt{heapq}, \texttt{collections}, \texttt{re}). The sandbox waits for the emulator after each button press so skills operate on fresh state. In our experiments, agents write BFS pathfinding routines that parse the ASCII map from game state and route around obstacles. In $\mathcal{H}_{\text{expert}}$ this is a pre-built A* tool; in AutoEvolve the agent writes it from scratch.

All stores support human-readable identifiers (e.g., \texttt{move\_to\_coords}), so the orchestrator and sub-agents reference capabilities by name.

\subsection{Memory}\label{sec:memory}

The memory store $\mathcal{M}$ is a persistent, hierarchically organized knowledge base. Each entry has a path (e.g., \texttt{mechanics/combat}), title, content, and importance score. The orchestrator's prompt includes a tree overview and can read entries by ID or title.

The agent writes memory during normal gameplay: game mechanics, locations, NPC info, team composition, strategies. The system prompt says to store what you learn but gives no game-specific guidance. The evolver acts as curator on top: filling gaps the agent missed, updating stale entries, lowering importance for areas the agent has moved past. Both sources are tracked with a \texttt{source} field and per-entry mutation history for post-hoc analysis.
